\chapter{多因子威胁度与差异化安全距离}

本章对应GTOC的第一项改进，针对Apollo PathBoundsDecider对所有障碍物使用统一安全距离的设计缺陷，提出基于多因子威胁度评估的差异化安全裕度$\delta_i$分配机制。本章首先深入剖析Apollo一刀切缓冲在多障碍物场景下压缩通行空间的源码根源与几何后果，进而设计融合TTC、横向重叠率、相对速度、障碍物类型与交互密度五个因子的威胁度评估模型，并将其映射为差异化的安全裕度。差异化$\delta_i$将作为基础量被第六章的最大间隙公式与第七章的时空走廊边界共同采用。

\section{安全裕度一刀切压缩通行空间}
\label{sec:problem_margin}

\subsection{问题描述}

PathBoundsDecider在构建路径边界时，需要在障碍物边缘外预留安全距离，以确保自车与障碍物之间保持足够的横向间距。该安全距离由\texttt{GetBuffer\-Between\-ADCCenter\-AndEdge}函数统一计算

\begin{lstlisting}[language=C++,caption={安全距离统一计算与nudge缓冲的静态/非静态分支},label={lst:buffer_calc}]
// path_bounds_decider_util.cc
// 相关 gflags 默认值（planning_gflags.cc）
//   obstacle_lat_buffer              = 0.4 m
//   static_obstacle_nudge_l_buffer   = 0.3 m
//   nonstatic_obstacle_nudge_l_buffer = 0.4 m
//   static_obstacle_speed_threshold  = 0.5 m/s

// 所有障碍物统一调用此函数计算安全距离
double PathBoundsDeciderUtil::
    GetBufferBetweenADCCenterAndEdge() {
  double adc_half_width =
      VehicleConfigHelper::GetConfig()
          .vehicle_param().width() / 2.0;
  return (adc_half_width + FLAGS_obstacle_lat_buffer);
}

// --- 静态/非静态的判定条件 ---
if (!obstacle.IsStatic() ||
    obstacle.speed() > FLAGS_static_obstacle_speed_threshold) {
  return false;  // 非静态障碍物不参与路径决策
}

// --- 边界松弛中的静态/非静态缓冲分支 ---
// 非静态障碍物使用 nonstatic buffer
double new_buffer = std::max(FLAGS_ego_front_slack_buffer,
                             FLAGS_nonstatic_obstacle_nudge_l_buffer);
// ...
// 静态障碍物使用 static buffer
if (path_bound_point->is_nudge_bound[LEFT_INDEX]) {
  old_buffer = std::max(FLAGS_obstacle_lat_buffer,
                        FLAGS_static_obstacle_nudge_l_buffer);
}
\end{lstlisting}

从代码可以看出，Apollo对安全距离的处理存在两个层面的粗粒度问题
\begin{enumerate}
\item \texttt{GetBufferBetweenADCCenterAndEdge}函数对所有障碍物返回相同的值，即半车宽加0.4\,m，完全不区分障碍物的类型、速度、横向偏移等特征。无论是远离车道的路锥、车道边缘的花坛，还是正对车道的违停车，均使用相同的安全距离。

\item Apollo虽然在\texttt{nudge\_l\_buffer}参数中区分了静态0.3\,m和非静态0.4\,m两个等级，但仅以0.5\,m/s的速度阈值做二分，低于该阈值视为静态，高于则视为非静态。该二分判定在 PathBoundsDecider 与 ObstacleNudgeDecider 的缓冲计算中通过 \texttt{!IsStatic() || speed() > FLAGS\_static\_obstacle\_speed\_threshold} 条件式生效。这种二分法无法区分运动方向随时突变的行人与轨迹可预测的同向行驶车辆，也无法区分碰撞时间短的正面接近障碍物与碰撞风险低的横向远离障碍物。
\end{enumerate}

\subsection{影响分析}

以下通过一个具体场景说明统一安全距离如何导致不必要的通行失败。定义净间距为道路边界到障碍物边缘之间扣除缓冲距离后的剩余宽度，自车可从某一侧通过当且仅当该侧净间距不小于自车宽度。

\textbf{场景设定}\quad 标准车道宽3.75\,m，$l_{road}^- = -1.875$\,m，$l_{road}^+ = 1.875$\,m。车道中偏右位置有一个路锥，横向占据$[-0.375,\; -0.075]$\,m，宽度0.3\,m。自车宽度$W_{ego} = 1.8$\,m。

\textbf{Apollo统一缓冲}\quad 取$d_{buf} = 0.3$\,m。从左侧通过时，净间距为
\[
1.875 - (-0.075) - 0.3 = 1.65\,\text{m} < W_{ego} = 1.8\,\text{m}
\]
从右侧通过时，净间距为
\[
(-0.375) - (-1.875) - 0.3 = 1.20\,\text{m} < 1.8\,\text{m}
\]
两侧均不足以容纳自车，路径被判定为BLOCKED，自车停车等待。

\textbf{差异化缓冲}\quad 该路锥处于静止状态且远离车道中心，威胁较低，取$d_{buf} = 0.15$\,m。从左侧通过时，净间距为
\[
1.875 - (-0.075) - 0.15 = 1.80\,\text{m} = W_{ego}
\]
恰好等于自车宽度，自车可从左侧安全通过。仅因缓冲参数从0.3\,m降至0.15\,m，净间距就从1.65\,m提升至1.80\,m，通行结果从BLOCKED翻转为可行。图\ref{fig:safety_margin_compression}直观对比了两种缓冲参数下的通行带差异。按障碍物威胁程度差异化分配安全距离的具体方法将在本章下一节展开。

\begin{figure}[htbp]
\centering
\resizebox{\textwidth}{!}{% 路锥场景：差异化 d_buf vs 统一 d_buf 对比
\begin{tikzpicture}[scale=1.0, transform shape]

% ============================================================
% 子图 (a)：差异化 d_buf=0.15m，恰好可通行
% ============================================================
\begin{scope}[shift={(0,0)}]
  % 背景
  \fill[bglight] (0,-2.3) rectangle (7.0,2.3);

  % 道路边界 ±1.875m
  \draw[deepblue, very thick] (0,-1.875) -- (7.0,-1.875);
  \draw[deepblue, very thick] (0, 1.875) -- (7.0, 1.875);
  \draw[deepblue, dashed, thin] (0,0) -- (7.0,0);

  % 车道宽度标注
  \draw[{Stealth}-{Stealth}, thin, deepblue] (0.3,-1.875) -- (0.3,1.875);
  \node[font=\scriptsize, deepblue, rotate=90] at (0.6,0) {车道 3.75m};

  % 路锥（宽0.3m，高0.3m，中心在 (3.25, -0.225)）
  % 横向 l: [-0.375, -0.075]，纵向 s: [3.1, 3.4]
  \fill[dangerred!70, draw=dangerred, thick] (3.1,-0.375) rectangle (3.4,-0.075);
  \node[white, font=\tiny\bfseries] at (3.25,-0.225) {锥};

  % 缓冲区 d_buf=0.15m（四周等距膨胀 0.15m）
  % l: [-0.375-0.15, -0.075+0.15] = [-0.525, 0.075]
  % s: [3.1-0.15, 3.4+0.15] = [2.95, 3.55]
  \draw[dangerred!70, dashed, thick] (2.95,-0.525) rectangle (3.55,0.075);
  \node[dangerred!70!black, font=\scriptsize, anchor=west] at (3.7,-0.225)
    {$\delta{=}0.15$m};

  % 膨胀距离标注（垂直）
  \draw[{Stealth}-{Stealth}, dangerred!60, thin] (2.8,-0.375) -- (2.8,-0.525);
  \node[dangerred!60, font=\tiny, left] at (2.75,-0.45) {0.15};
  \draw[{Stealth}-{Stealth}, dangerred!60, thin] (2.8,-0.075) -- (2.8,0.075);
  \node[dangerred!60, font=\tiny, left] at (2.75,0.0) {0.15};

  % 通行区域 [0.075, 1.875]（绿色填充）
  \fill[lightgreen, opacity=0.25] (1.0,0.075) rectangle (5.5,1.875);

  % 自车（宽1.8m = y方向，长4.5m = x方向，1:1比例）
  \fill[deepblue, opacity=0.3] (1.0,0.075) rectangle (5.5,1.875);
  \draw[deepblue, thick, rounded corners=2pt] (1.0,0.075) rectangle (5.5,1.875);
  \node[deepblue, font=\scriptsize\bfseries] at (3.25,0.975) {自车};

  % 净间距标注
  \draw[{Stealth}-{Stealth}, thick, lightgreen!60!black] (6.2,0.075) -- (6.2,1.875);
  \node[font=\scriptsize, lightgreen!60!black, anchor=west] at (6.3,0.975) {1.80m};
  \node[font=\scriptsize, lightgreen!60!black, anchor=west] at (6.3,0.55) {$= W_{\mathrm{ego}}$};

  \node[lightgreen!60!black, font=\normalsize\bfseries] at (1.5,-2.1) {\checkmark};

  % 子图标题
  \node[font=\small\bfseries, align=center] at (3.5,-2.7)
    {(a) $\delta{=}0.15$m（差异化），恰好可通行};
\end{scope}

% ============================================================
% 子图 (b)：Apollo d_buf=0.3m，BLOCKED
% ============================================================
\begin{scope}[shift={(8.5,0)}]
  % 背景
  \fill[bglight] (0,-2.3) rectangle (7.0,2.3);

  % 道路边界 ±1.875m
  \draw[deepblue, very thick] (0,-1.875) -- (7.0,-1.875);
  \draw[deepblue, very thick] (0, 1.875) -- (7.0, 1.875);
  \draw[deepblue, dashed, thin] (0,0) -- (7.0,0);

  % 车道宽度标注
  \draw[{Stealth}-{Stealth}, thin, deepblue] (0.3,-1.875) -- (0.3,1.875);
  \node[font=\scriptsize, deepblue, rotate=90] at (0.6,0) {车道 3.75m};

  % 路锥（同样大小）
  \fill[dangerred!70, draw=dangerred, thick] (3.1,-0.375) rectangle (3.4,-0.075);
  \node[white, font=\tiny\bfseries] at (3.25,-0.225) {锥};

  % 缓冲区 d_buf=0.3m（四周等距膨胀 0.3m）
  % l: [-0.375-0.3, -0.075+0.3] = [-0.675, 0.225]
  % s: [3.1-0.3, 3.4+0.3] = [2.8, 3.7]
  \draw[dangerred!70, dashed, thick] (2.8,-0.675) rectangle (3.7,0.225);
  \node[dangerred!70!black, font=\scriptsize, anchor=west] at (3.85,-0.225)
    {$\delta{=}0.3$m};

  % 膨胀距离标注
  \draw[{Stealth}-{Stealth}, dangerred!60, thin] (2.6,-0.375) -- (2.6,-0.675);
  \node[dangerred!60, font=\tiny, left] at (2.55,-0.525) {0.3};
  \draw[{Stealth}-{Stealth}, dangerred!60, thin] (2.6,-0.075) -- (2.6,0.225);
  \node[dangerred!60, font=\tiny, left] at (2.55,0.075) {0.3};

  % 上方不足区域 [0.225, 1.875] = 1.65m
  \fill[dangerred, opacity=0.12] (1.0,0.225) rectangle (5.5,1.875);
  % 下方不足区域 [-1.875, -0.675] = 1.2m
  \fill[dangerred, opacity=0.12] (1.0,-1.875) rectangle (5.5,-0.675);

  % 上方净间距
  \draw[{Stealth}-{Stealth}, thick, dangerred] (5.7,0.225) -- (5.7,1.875);
  \node[font=\scriptsize, dangerred, anchor=west] at (5.8,1.05) {1.65m};
  \node[font=\scriptsize, dangerred, anchor=west] at (5.8,0.65) {$< W_{\mathrm{ego}}$};

  % 下方净间距
  \draw[{Stealth}-{Stealth}, thick, dangerred] (5.7,-1.875) -- (5.7,-0.675);
  \node[font=\scriptsize, dangerred, anchor=west] at (5.8,-1.275) {1.20m};
  \node[font=\scriptsize, dangerred, anchor=west] at (5.8,-1.6) {$< W_{\mathrm{ego}}$};

  % BLOCKED
  \draw[dangerred, ultra thick] (2.5,1.0) -- (4.0,-1.0);
  \draw[dangerred, ultra thick] (2.5,-1.0) -- (4.0,1.0);
  \node[dangerred, font=\small\bfseries] at (3.25,-2.1) {BLOCKED};

  % 子图标题
  \node[font=\small\bfseries, align=center] at (3.5,-2.7)
    {(b) $\delta{=}0.3$m（统一），BLOCKED};
\end{scope}

\end{tikzpicture}
}
\caption{统一缓冲与差异化缓冲下的通行带对比}
\label{fig:safety_margin_compression}
\end{figure}
\section{多因子威胁度评估与动态安全距离}
上一节识别了Baseline在安全裕度维度的次优性，式\eqref{eq:safety_margin}中的统一$\delta$不区分障碍物威胁程度，导致低威胁障碍物过度压缩通行带。本节建立量化模型，为每个障碍物计算综合威胁度$\Theta_i$，并据此确定差异化的安全裕度$\delta_i$。

\subsection{威胁度计算模型}
对第$i$个障碍物，其综合威胁度$\Theta_i \in [0, 1]$由五个因子加权求和得到

\begin{equation}
    \Theta_i = w_1 \cdot f_{TTC}(i) + \\
    w_2 \cdot f_{overlap}(i) + \\
    w_3 \cdot f_{vel}(i) + \\
    w_4 \cdot f_{type}(i) + \\
    w_5 \cdot f_{inter}(i)
    \label{eq:threat_score}
\end{equation}

权重默认取值为$w_1 = 0.30$，$w_2 = 0.20$，$w_3 = 0.15$，$w_4 = 0.10$，$w_5 = 0.25$。上述权重通过以下方式确定。首先基于各因子对多障碍物场景区分度的定性分析确定初始权重，然后通过对称路锥、横穿行人等典型场景的定性观察调整至最终取值。TTC因子和交互密度因子对场景风险的区分度最高，因此赋予较大权重，障碍物类型在同一场景中变化较小，赋予较小权重。各因子的取值范围均归一化到$[0, 1]$，定义如下。

\subsubsection{碰撞时间因子$f_{TTC}$}
碰撞时间（Time-to-Collision, TTC）是衡量碰撞紧迫性的经典指标

\begin{equation}
    TTC_i = \frac{s_i - s_{ego}}{\dot{s}_{ego} - \dot{s}_i}
    \label{eq:ttc}
\end{equation}

其中$s_i$为障碍物纵向位置，$\dot{s}$为纵向速度。当$\dot{s}_{ego} \leq \dot{s}_i$即自车不会追上障碍物时，$TTC_i = +\infty$。TTC因子通过分段线性函数映射至$[0, 1]$

\begin{equation}
    f_{TTC}(i) = \begin{cases}
        1.0 & TTC_i \leq T_{crit} \\[0.3em]
        \displaystyle\frac{T_{max} - TTC_i}{T_{max} - T_{crit}} & T_{crit} < TTC_i < T_{max} \\[0.5em]
        0 & TTC_i \geq T_{max}
    \end{cases}
    \label{eq:f_ttc}
\end{equation}

其中$T_{crit} = 2$\,s为临界碰撞时间，$T_{max} = 7$\,s为最大关注时间，与Apollo ST图的时间范围一致。$TTC_i \leq 2$\,s意味着碰撞迫在眉睫，威胁最大；$TTC_i \geq 7$\,s的障碍物已超出当前规划周期的关注范围。

\subsubsection{横向重叠率因子$f_{overlap}$}
横向重叠率衡量障碍物与自车行驶走廊在横向上的重叠程度

\begin{equation}
    f_{overlap}(i) = \frac{\max\!\left(0,\; \min(l_i^{+}, l_{ego}^{+}) - \max(l_i^{-}, l_{ego}^{-})\right)}{W_{ego}}
    \label{eq:f_overlap}
\end{equation}

其中$l_{ego}^{+} = l_{ego} + W_{ego}/2$、$l_{ego}^{-} = l_{ego} - W_{ego}/2$为自车走廊的横向边界。$f_{overlap} = 0$表示完全不重叠即障碍物在自车投影之外，$f_{overlap} = 1$表示完全重叠即障碍物正对自车。该因子直接反映了碰撞的“几何可能性”，横向偏移较大的障碍物即使TTC较小，实际碰撞风险也低。

\subsubsection{相对速度因子$f_{vel}$}
相对速度因子反映障碍物与自车的接近速度

\begin{equation}
    f_{vel}(i) = \sigma\!\left(\frac{v_{rel,i}}{v_{max}}\right), \quad \sigma(x) = \frac{1}{1 + e^{-5x}}
    \label{eq:f_vel}
\end{equation}

其中$v_{rel,i} = \dot{s}_{ego} - \dot{s}_i$为相对纵向速度，正值表示自车接近障碍物，$v_{max}$为道路限速。Sigmoid函数将相对速度平滑映射至$(0, 1)$，快速接近的障碍物即$v_{rel} \gg 0$时趋近1，同向远离的障碍物即$v_{rel} < 0$时趋近0。选择Sigmoid而非线性映射的原因在于，相对速度为零附近的区分度更重要，即低速接近和低速远离的差异，而极高相对速度间的差异不需要精细区分。

\subsubsection{障碍物类型因子$f_{type}$}
\textbf{取值映射}\quad{}不同类型障碍物的运动不确定性与受伤严重度差异显著，按下式取离散值

\begin{equation}
    f_{type}(i) = \begin{cases}
        1.0 & \text{行人、自行车等弱势道路使用者} \\
        0.7 & \text{机动车，运动相对可预测} \\
        0.5 & \text{未知可动目标，按保守权重处理} \\
        0.3 & \text{静态障碍物，仅需基本缓冲}
    \end{cases}
    \label{eq:f_type}
\end{equation}

\textbf{取值依据}\quad{}该映射由运动不确定性与事故后果两方面共同确定。一方面，行人和自行车的运动方向和速度随时可能突变，如突然横穿或急停回转，其轨迹可预测性远低于机动车，需要最大安全裕度；机动车受道路与交通规则约束，轨迹相对可预测；静态障碍物无运动不确定性，仅保留基本缓冲。另一方面，交通事故统计表明弱势道路使用者（Vulnerable Road Users, VRU）在同等碰撞能量下受伤严重度和致死率远高于车内乘员，规划器应对其施加额外保护性裕度，因此赋予行人和自行车最高权重1.0。

\textbf{Apollo感知映射}\quad{}Apollo感知模块在\texttt{PerceptionObstacle}消息中提供枚举字段\texttt{type}，取值覆盖\texttt{UNKNOWN}、\texttt{UNKNOWN\_MOVABLE}、\texttt{UNKNOWN\_UNMOVABLE}、\texttt{PEDESTRIAN}、\texttt{BICYCLE}和\texttt{VEHICLE}六类。规划层直接查表即可完成从感知标签到$f_{type}$的映射，\texttt{PEDESTRIAN}和\texttt{BICYCLE}归入1.0分支，\texttt{VEHICLE}归入0.7分支，\texttt{UNKNOWN\_MOVABLE}与\texttt{UNKNOWN}归入0.5分支，\texttt{UNKNOWN\_UNMOVABLE}归入0.3分支。对未知可动目标采用中等偏上的保守权重是因为分类置信度不足时宁可过估威胁也不可低估，避免将潜在行人误判为低威胁物体。

\textbf{合成方式}\quad{}$f_{type}$作为乘性类别因子进入式\eqref{eq:threat_score}，与碰撞时间、横向重叠率、相对速度三项连续因子加权叠加，合成综合威胁度$\Theta_i$，并经式\eqref{eq:dynamic_delta}映射为差异化安全裕度$\delta_i$，实现对VRU的主动倾斜保护。

\subsubsection{交互密度因子$f_{inter}$}
交互密度因子量化障碍物$i$与周围其他障碍物的聚集程度，是本文威胁度模型中与分组策略直接相关的因子

\begin{equation}
    f_{inter}(i) = \frac{1}{N-1} \sum_{j \neq i} \mathbb{1}\!\left[d_{ij} < d_{cluster}\right] \cdot \frac{d_{cluster} - d_{ij}}{d_{cluster}}
    \label{eq:f_inter}
\end{equation}

其中$N$为障碍物总数，$d_{ij} = \sqrt{(s_i - s_j)^2 + (l_i - l_j)^2}$为障碍物$i$和$j$在SL空间中的欧氏距离，$d_{cluster} = 10$\,m为聚类半径，$\mathbb{1}[\cdot]$为指示函数。
$f_{inter}$的物理含义是，若障碍物$i$周围有大量其他障碍物聚集，则该因子值高，表明局部交互密度大，需要更谨慎的避让裕度。孤立障碍物的$f_{inter} \approx 0$，密集车群中的障碍物$f_{inter}$接近1。

引入交互密度因子的意义在于它能区分“3个分散的障碍物”和“3个聚集的障碍物”这两种对避让策略要求截然不同的场景，前者各自独立，安全裕度可适当放宽，后者构成密集障碍物组，留给自车的横向空间本就有限，需要更精确的安全边界。这一因子与下一节的纵向重叠分组形成自然呼应，$f_{inter}$高的障碍物大概率属于同一连通分量，在分组内被统一协调。

图\ref{fig:threat_factors}以子图形式给出了各因子的函数形态。$f_\text{TTC}$在$T_\text{crit}$以下取饱和值1，在$T_\text{max}$以上回到0，中间为线性过渡；$f_\text{overlap}$随重叠率线性增长；$f_\text{vel}$采用Sigmoid函数，对零相对速度附近的变化最为敏感；$f_\text{type}$为离散取值，反映不同障碍物类型的固有风险差异；$f_\text{inter}$中单个邻居的贡献随距离线性衰减至聚类半径$d_{cluster}$处归零。各因子均归一化到$[0,1]$，权重之和为1，综合威胁度$\Theta_i$也落在$[0,1]$区间，可直接作为式\eqref{eq:dynamic_delta}中安全裕度线性插值的参数。

\begin{figure}[htbp]
\centering
% 威胁度五因子函数曲线
% 每个 tikzpicture 用 \useasboundingbox 统一尺寸，\resizebox 缩放不变形

\begin{subfigure}[b]{0.32\textwidth}
\centering
\resizebox{\textwidth}{!}{%
\begin{tikzpicture}
  \useasboundingbox (-0.5,-0.8) rectangle (3.2,3.0);
  \draw[->, thick] (0,0) -- (2.8,0) node[right,font=\scriptsize] {TTC (s)};
  \draw[->, thick] (0,0) -- (0,2.8) node[above,font=\scriptsize] {$f_{\text{TTC}}$};
  \node[font=\scriptsize, left] at (0,2.4) {1.0};
  \draw[gray, thin] (-0.08,2.4) -- (0.08,2.4);
  \node[font=\scriptsize, below] at (0.6,0) {2};
  \node[font=\scriptsize, below] at (2.1,0) {7};
  \draw[gray, thin] (0.6,-0.08) -- (0.6,0.08);
  \draw[gray, thin] (2.1,-0.08) -- (2.1,0.08);
  \draw[deeppink, very thick] (0,2.4) -- (0.6,2.4) -- (2.1,0) -- (2.6,0);
  \draw[gray, dashed, thin] (0.6,0) -- (0.6,2.4);
  \node[font=\scriptsize] at (0.6,2.65) {$T_{crit}$};
  \node[font=\scriptsize] at (2.1,0.25) {$T_{max}$};
\end{tikzpicture}%
}
\caption{碰撞时间因子}
\end{subfigure}
\hfill
\begin{subfigure}[b]{0.32\textwidth}
\centering
\resizebox{\textwidth}{!}{%
\begin{tikzpicture}
  \useasboundingbox (-0.5,-0.8) rectangle (3.2,3.0);
  \draw[->, thick] (0,0) -- (2.8,0) node[right,font=\scriptsize] {重叠率};
  \draw[->, thick] (0,0) -- (0,2.8) node[above,font=\scriptsize] {$f_{\text{ov}}$};
  \node[font=\scriptsize, left] at (0,2.4) {1.0};
  \draw[gray, thin] (-0.08,2.4) -- (0.08,2.4);
  \node[font=\scriptsize, below] at (2.4,0) {1.0};
  \node[font=\scriptsize, below] at (0,0) {0};
  \draw[deeppink, very thick] (0,0) -- (2.4,2.4);
\end{tikzpicture}%
}
\caption{横向重叠率因子}
\end{subfigure}
\hfill
\begin{subfigure}[b]{0.32\textwidth}
\centering
\resizebox{\textwidth}{!}{%
\begin{tikzpicture}
  \useasboundingbox (-0.5,-0.8) rectangle (3.2,3.0);
  % 内容居中: sigmoid 画在 [0.4, 2.8] 范围, 对称轴在 1.6
  \draw[->, thick] (0.2,0) -- (3.0,0) node[right,font=\scriptsize] {$v_{rel}/v_{max}$};
  \draw[->, thick] (0.2,0) -- (0.2,2.8) node[above,font=\scriptsize] {$f_{vel}$};
  \node[font=\scriptsize, left] at (0.2,2.4) {1.0};
  \draw[gray, thin] (0.12,2.4) -- (0.28,2.4);
  \node[font=\scriptsize, left] at (0.2,1.2) {0.5};
  \draw[gray, thin] (0.12,1.2) -- (0.28,1.2);
  % v_rel/v_max: -1 → x=0.4, 0 → x=1.6, 1 → x=2.8
  \draw[deeppink, very thick]
    plot[smooth, domain=0.4:2.8, samples=50] (\x, {2.4/(1+exp(-5*(\x-1.6)/1.2))});
  \node[font=\scriptsize, below] at (0.4,0) {$-1$};
  \node[font=\scriptsize, below] at (1.6,0) {$0$};
  \node[font=\scriptsize, below] at (2.8,0) {$1$};
  \draw[gray, dashed, thin] (1.6,0) -- (1.6,1.2);
\end{tikzpicture}%
}
\caption{相对速度因子}
\end{subfigure}

\vspace{0.8cm}

\hfill
\begin{subfigure}[b]{0.32\textwidth}
\centering
\resizebox{\textwidth}{!}{%
\begin{tikzpicture}
  \useasboundingbox (-0.5,-0.8) rectangle (3.2,3.0);
  \draw[->, thick] (0,0) -- (2.8,0);
  \draw[->, thick] (0,0) -- (0,2.8) node[above,font=\scriptsize] {$f_{\text{type}}$};
  \node[font=\scriptsize, left] at (0,2.4) {1.0};
  \draw[gray, thin] (-0.08,2.4) -- (0.08,2.4);
  \node[font=\scriptsize, left] at (0,1.68) {0.7};
  \draw[gray, thin] (-0.08,1.68) -- (0.08,1.68);
  \node[font=\scriptsize, left] at (0,0.72) {0.3};
  \draw[gray, thin] (-0.08,0.72) -- (0.08,0.72);
  \fill[dangerred!60] (0.2,0) rectangle (0.7,2.4);
  \fill[lightorange!80] (0.9,0) rectangle (1.4,1.68);
  \fill[deepblue!40] (1.6,0) rectangle (2.1,0.72);
  \node[font=\scriptsize] at (0.45,-0.35) {行人};
  \node[font=\scriptsize] at (1.15,-0.35) {车辆};
  \node[font=\scriptsize] at (1.85,-0.35) {静态};
\end{tikzpicture}%
}
\caption{障碍物类型因子}
\end{subfigure}
\hspace{0.02\textwidth}
\begin{subfigure}[b]{0.32\textwidth}
\centering
\resizebox{\textwidth}{!}{%
\begin{tikzpicture}
  \useasboundingbox (-0.5,-0.8) rectangle (3.2,3.0);
  \draw[->, thick] (0,0) -- (2.8,0) node[right,font=\scriptsize] {$d_{ij}$ (m)};
  \draw[->, thick] (0,0) -- (0,2.8) node[above,font=\scriptsize] {贡献};
  \node[font=\scriptsize, left] at (0,2.4) {1.0};
  \draw[gray, thin] (-0.08,2.4) -- (0.08,2.4);
  \node[font=\scriptsize, below] at (2.4,0) {$d_{cl}{=}10$};
  \draw[gray, thin] (2.4,-0.08) -- (2.4,0.08);
  \draw[deeppink, very thick] (0,2.4) -- (2.4,0) -- (2.7,0);
  \node[font=\scriptsize, deeppink] at (1.6,2.0) {$\frac{d_{cl}-d_{ij}}{d_{cl}}$};
\end{tikzpicture}%
}
\caption{单邻居交互贡献}
\end{subfigure}
\hfill

\caption{多因子威胁度评估的五个因子函数形态}
\label{fig:threat_factors}
\end{figure}

图\ref{fig:threat_radar}进一步给出了两个典型场景下的五维威胁度雷达图对照。左侧为对向来车场景，TTC较小、横向重叠率高、相对速度大、对象为车辆、交互密度中等，五维雷达图呈现大面积覆盖，综合威胁度接近0.85，对应$\delta_i$取到$\delta_\text{max}$附近，安全裕度最大。右侧为静止路锥场景，TTC趋于无穷、重叠率低、相对速度为零、类型权重低、密度因子接近0，雷达图呈现小面积覆盖，综合威胁度不足0.2，对应$\delta_i$接近$\delta_\text{min}$，从而大幅释放横向空间。这两张雷达图直观说明了差异化安全裕度的判别逻辑与最终数值效应。

\begin{figure}[htbp]
\centering
\begin{subfigure}[b]{0.48\textwidth}
  \centering
  \resizebox{\linewidth}{!}{% 高威胁障碍物雷达图（子图 a）
\begin{tikzpicture}[scale=0.85, transform shape]
  \foreach \r in {1, 2, 3} {
    \draw[gray!40, thin] (90:\r) \foreach \a in {162, 234, 306, 378} { -- (\a:\r) } -- cycle;
  }
  \foreach \a/\lbl in {90/$f_{TTC}$, 162/$f_{overlap}$, 234/$f_{vel}$, 306/$f_{type}$, 378/$f_{inter}$} {
    \draw[gray!60, thin] (0,0) -- (\a:3.3);
    \node[font=\footnotesize] at (\a:3.7) {\lbl};
  }

  \fill[dangerred!30, draw=dangerred, thick]
    (90:2.7) -- (162:2.4) -- (234:2.55) -- (306:2.1) -- (378:1.8) -- cycle;
  \foreach \a/\r in {90/2.7, 162/2.4, 234/2.55, 306/2.1, 378/1.8} {
    \fill[dangerred] (\a:\r) circle (2.5pt);
  }

  \node[draw=dangerred, fill=dangerred!15, rounded corners, font=\small\bfseries] at (0, -4) {$\Theta = 0.82$ (CRITICAL)};

  \node[gray, font=\tiny] at (-3.8, 0) {0};
  \node[gray, font=\tiny] at (-3.8, 1.0) {0.33};
  \node[gray, font=\tiny] at (-3.8, 2.0) {0.67};
  \node[gray, font=\tiny] at (-3.8, 3.0) {1.0};
\end{tikzpicture}
}
  \caption{高威胁 \;\; 对向来车}
  \label{fig:threat_radar:a}
\end{subfigure}\hfill
\begin{subfigure}[b]{0.48\textwidth}
  \centering
  \resizebox{\linewidth}{!}{% 低威胁障碍物雷达图（子图 b）
\begin{tikzpicture}[scale=0.85, transform shape]
  \foreach \r in {1, 2, 3} {
    \draw[gray!40, thin] (90:\r) \foreach \a in {162, 234, 306, 378} { -- (\a:\r) } -- cycle;
  }
  \foreach \a/\lbl in {90/$f_{TTC}$, 162/$f_{overlap}$, 234/$f_{vel}$, 306/$f_{type}$, 378/$f_{inter}$} {
    \draw[gray!60, thin] (0,0) -- (\a:3.3);
    \node[font=\footnotesize] at (\a:3.7) {\lbl};
  }

  \fill[lightgreen!40, draw=lightgreen!70!black, thick]
    (90:0.3) -- (162:0.6) -- (234:0.3) -- (306:0.9) -- (378:0.45) -- cycle;
  \foreach \a/\r in {90/0.3, 162/0.6, 234/0.3, 306/0.9, 378/0.45} {
    \fill[lightgreen!70!black] (\a:\r) circle (2.5pt);
  }

  \node[draw=lightgreen!70!black, fill=lightgreen!15, rounded corners, font=\small\bfseries] at (0, -4) {$\Theta = 0.18$ (MINOR)};
\end{tikzpicture}
}
  \caption{低威胁 \;\; 远处静态路锥}
  \label{fig:threat_radar:b}
\end{subfigure}
\caption{典型场景下的威胁度雷达图对比}
\label{fig:threat_radar}
\end{figure}

\subsection{动态安全距离}
基于综合威胁度$\Theta_i$，将式\eqref{eq:safety_margin}中的统一安全裕度替换为差异化版本

\begin{equation}
    \delta_i = \delta_{min} + (\delta_{max} - \delta_{min}) \cdot \Theta_i
    \label{eq:dynamic_delta}
\end{equation}

其中$\delta_{max} = W_{ego}/2 + d_{buf,max}$为最大安全裕度，对应CRITICAL等级障碍物，$\delta_{min} = W_{ego}/2 + d_{buf,min}$为最小安全裕度，对应MINOR等级障碍物。$d_{buf,max}$与Apollo原始\texttt{point\_extension}取值一致，$d_{buf,min}$取较小值以释放低威胁障碍物占用的横向空间，两者的具体取值由系统集成阶段的参数标定确定。

式\eqref{eq:dynamic_delta}的设计确保了以下性质。CRITICAL障碍物时$\Theta_i \to 1$，获得与Apollo原始方案相同或更大的安全裕度，不降低安全性；MINOR障碍物时$\Theta_i \to 0$，获得较小但仍安全的裕度，释放被不必要占用的横向空间。

相应地，式\eqref{eq:def_u}--\eqref{eq:def_v}中的关键量更新为差异化版本

\begin{align}
    u_i &= l_i^{-} - \delta_i \label{eq:def_u_dynamic} \\
    v_i &= l_i^{+} + \delta_i \label{eq:def_v_dynamic}
\end{align}

当$\Theta_i$对所有障碍物均取相同值时，式\eqref{eq:def_u_dynamic}--\eqref{eq:def_v_dynamic}退化为式\eqref{eq:def_u}--\eqref{eq:def_v}。动态$\delta_i$的引入不影响后续引理1和定理1的成立，两者的证明仅依赖$u$和$v$值的排序关系，只需按$u_i$而非$l_i^{-}$排序即可。

\section{本章小结}

本章针对Apollo PathBoundsDecider安全裕度一刀切的设计缺陷，提出GTOC的第一项改进，即多因子威胁度评估与差异化安全距离机制。在源码层面定位了\texttt{GetBuffer\-Between\-ADCCenter\-AndEdge}函数对所有障碍物返回相同安全距离的根源，并通过几何分析量化了统一缓冲对通行带宽度的不必要压缩。在算法设计层面，构建了融合TTC、横向重叠率、相对速度、障碍物类型与交互密度的多因子威胁度模型，并将归一化威胁度通过分段映射转化为差异化的安全裕度$\delta_i$。差异化$\delta_i$不依赖统一缓冲假设，可直接代入第六章的最大间隙公式与第七章的时空走廊边界构造，作为后续两章设计的共用输入。
